\documentclass[11pt,fleqn, openany]{book} % Default font size and left-justified equations

%%%%%%%%%%%%%%%%%%%%%%%%%%%%%%%%%%%%%%%%%
% The Legrand Orange Book
% Structural Definitions File
% Version 2.1 (26/09/2018)
%
% Original author:
% Mathias Legrand (legrand.mathias@gmail.com) with modifications by:
% Vel (vel@latextemplates.com)
% 
% This file was downloaded from:
% http://www.LaTeXTemplates.com
%
% License:
% CC BY-NC-SA 3.0 (http://creativecommons.org/licenses/by-nc-sa/3.0/)
%
%%%%%%%%%%%%%%%%%%%%%%%%%%%%%%%%%%%%%%%%%

%----------------------------------------------------------------------------------------
%	VARIOUS REQUIRED PACKAGES AND CONFIGURATIONS
%----------------------------------------------------------------------------------------

\usepackage[table]{xcolor}

\usepackage{graphicx}
\usepackage{tabularx} % Required for including pictures
\usepackage{pgf,tikz,tkz-tab,eurosym,yhmath, stmaryrd}
\usepackage{pgfplots}
\usepackage{mathrsfs}
\usetikzlibrary{patterns}
\usetikzlibrary{trees}
\graphicspath{{../../Pictures/}}
\usepackage{multicol} 


\usepackage[english]{babel} % English language/hyphenation
\usepackage{icomma}
\usepackage{enumitem} % Customize lists
\setlist{nolistsep, nosep, nolistsep} % Reduce spacing between bullet points and numbered lists

\usepackage{booktabs} % Required for nicer horizontal rules in tables

 % Required for specifying colors by name


\definecolor{ocre}{RGB}{243,102,25} % Define the orange color used for highlighting throughout the book

\usepackage{listings}

\definecolor{codegreen}{rgb}{0,0.6,0}
\definecolor{codegray}{rgb}{0.5,0.5,0.5}
\definecolor{codepurple}{rgb}{0.58,0,0.82}
\definecolor{backcolour}{rgb}{0.95,0.95,0.92}

\lstdefinestyle{mystyle}{
    backgroundcolor=\color{backcolour},   
    commentstyle=\color{codegreen},
    keywordstyle=\color{magenta},
    numberstyle=\tiny\color{codegray},
    stringstyle=\color{codepurple},
    basicstyle=\ttfamily\footnotesize,
    breakatwhitespace=false,         
    breaklines=true,                 
    captionpos=b,                    
    keepspaces=true,                 
    numbers=left,                    
    numbersep=5pt,                  
    showspaces=false,                
    showstringspaces=false,
    showtabs=false,                  
    tabsize=2
}

\lstset{style=mystyle}

%----------------------------------------------------------------------------------------
% Paramétrage XSIM
%----------------------------------------------------------------------------------------

\usepackage[no-files]{xsim}


\DeclareExerciseEnvironmentTemplate{myex}{%
    \textbf{%
      \hypertarget{ex:\ExerciseID}{\sffamily{\ensuremath{\blacktriangleright}} Exercice \GetExerciseProperty{counter} \GetExerciseProperty{subtitle} --}
      \hyperlink{sol:\ExerciseID}{Voir le corrigé}%
    }\par
}{\par\smallskip}

\DeclareExerciseEnvironmentTemplate{mysol}{%
    \textbf{%
      \hypertarget{sol:\ExerciseID}{\sffamily{\ensuremath{\blacktriangleright}} Correction \GetExerciseProperty{counter} --}
      \hyperlink{ex:\ExerciseID}{Voir l'énoncé}%
    }\par
}{\par\medskip}

\xsimsetup{
  exercise/template = myex ,
  solution/template = mysol 
}

%Collection exercices

\DeclareExerciseTagging{topic}

\xsimsetup{collect}

%----------------------------------------------------------------------------------------
% SYMBOLES
%----------------------------------------------------------------------------------------

\newcommand\imCMsym[4][\mathord]{%
  \DeclareFontFamily{U} {#2}{}
  \DeclareFontShape{U}{#2}{m}{n}{
    <-6> #25
    <6-7> #26
    <7-8> #27
    <8-9> #28
    <9-10> #29
    <10-12> #210
    <12-> #212}{}
  \DeclareSymbolFont{CM#2} {U} {#2}{m}{n}
  \DeclareMathSymbol{#4}{#1}{CM#2}{#3}
}
\newcommand\alsoimCMsym[4][\mathord]{\DeclareMathSymbol{#4}{#1}{CM#2}{#3}}

\imCMsym{cmmi}{124}{\CMjmath}

\newcommand{\Oij}{(O\,;\,\vec{\imath}\,,\, \vec{\CMjmath} )}
\newcommand{\Oijk}{(O\,;\,\vec{\imath}\,,\, \vec{\CMjmath}\,,\,\vec{k})}

\newcommand\e{\mathrm{e}}
\newcommand\R{\mathbb{R}}
\newcommand\N{\mathbb{N}}


%----------------------------------------------------------------------------------------
%	MARGINS
%----------------------------------------------------------------------------------------

\usepackage{geometry} % Required for adjusting page dimensions and margins

\geometry{
	paper=a4paper, % Paper size, change to letterpaper for US letter size
	top=3cm, % Top margin
	bottom=3cm, % Bottom margin
	left=2cm, % Left margin
	right=2cm, % Right margin
	headheight=14pt, % Header height
	footskip=1.4cm, % Space from the bottom margin to the baseline of the footer
	headsep=10pt, % Space from the top margin to the baseline of the header
	%showframe, % Uncomment to show how the type block is set on the page
}

\setlength{\parindent}{0pt}
\parskip=5pt



%----------------------------------------------------------------------------------------
%	FONTS
%----------------------------------------------------------------------------------------

\usepackage{avant} % Use the Avantgarde font for headings
\usepackage{times} % Use the Times font for headings
\usepackage{mathptmx} % Use the Adobe Times Roman as the default text font together with math symbols from the Sym­bol, Chancery and Com­puter Modern fonts

%\usepackage{microtype} % Slightly tweak font spacing for aesthetics
%\usepackage[utf8]{inputenc} % Required for including letters with accents
\usepackage[T1]{fontenc} % Use 8-bit encoding that has 256 glyphs

%----------------------------------------------------------------------------------------
%	BIBLIOGRAPHY AND INDEX
%----------------------------------------------------------------------------------------

\usepackage[style=numeric,citestyle=numeric,sorting=nyt,sortcites=true,autopunct=true,babel=hyphen,hyperref=true,abbreviate=false,backref=true,backend=biber]{biblatex}
\addbibresource{bibliography.bib} % BibTeX bibliography file
\defbibheading{bibempty}{}

\usepackage{calc} % For simpler calculation - used for spacing the index letter headings correctly
\usepackage{makeidx} % Required to make an index
\makeindex % Tells LaTeX to create the files required for indexing

%----------------------------------------------------------------------------------------
%	MAIN TABLE OF CONTENTS
%----------------------------------------------------------------------------------------

\usepackage{titletoc} % Required for manipulating the table of contents

\contentsmargin{0cm} % Removes the default margin

% Part text styling (this is mostly taken care of in the PART HEADINGS section of this file)
\titlecontents{part}
	[0cm] % Left indentation
	{\addvspace{20pt}\bfseries} % Spacing and font options for parts
	{}
	{}
	{}

% Chapter text styling
\titlecontents{chapter}
	[1.25cm] % Left indentation
	{\addvspace{12pt}\large\sffamily\bfseries} % Spacing and font options for chapters
	{\color{ocre!60}\contentslabel[\Large\thecontentslabel]{1.25cm}\color{ocre}} % Formatting of numbered sections of this type
	{\color{ocre}} % Formatting of numberless sections of this type
	{\color{ocre!60}\normalsize\;\titlerule*[.5pc]{.}\;\thecontentspage} % Formatting of the filler to the right of the heading and the page number

% Section text styling
\titlecontents{section}
	[1.25cm] % Left indentation
	{\addvspace{3pt}\sffamily\bfseries} % Spacing and font options for sections
	{\contentslabel[\thecontentslabel]{1.25cm}} % Formatting of numbered sections of this type
	{} % Formatting of numberless sections of this type
	{\hfill\color{black}\thecontentspage} % Formatting of the filler to the right of the heading and the page number

% Subsection text styling
\titlecontents{subsection}
	[1.25cm] % Left indentation
	{\addvspace{1pt}\sffamily\small} % Spacing and font options for subsections
	{\contentslabel[\thecontentslabel]{1.25cm}} % Formatting of numbered sections of this type
	{} % Formatting of numberless sections of this type
	{\ \titlerule*[.5pc]{.}\;\thecontentspage} % Formatting of the filler to the right of the heading and the page number

% Figure text styling
\titlecontents{figure}
	[1.25cm] % Left indentation
	{\addvspace{1pt}\sffamily\small} % Spacing and font options for figures
	{\thecontentslabel\hspace*{1em}} % Formatting of numbered sections of this type
	{} % Formatting of numberless sections of this type
	{\ \titlerule*[.5pc]{.}\;\thecontentspage} % Formatting of the filler to the right of the heading and the page number

% Table text styling
\titlecontents{table}
	[1.25cm] % Left indentation
	{\addvspace{1pt}\sffamily\small} % Spacing and font options for tables
	{\thecontentslabel\hspace*{1em}} % Formatting of numbered sections of this type
	{} % Formatting of numberless sections of this type
	{\ \titlerule*[.5pc]{.}\;\thecontentspage} % Formatting of the filler to the right of the heading and the page number

%----------------------------------------------------------------------------------------
%	MINI TABLE OF CONTENTS IN PART HEADS
%----------------------------------------------------------------------------------------

% Chapter text styling
\titlecontents{lchapter}
	[0em] % Left indentation
	{\addvspace{15pt}\large\sffamily\bfseries} % Spacing and font options for chapters
	{\color{ocre}\contentslabel[\Large\thecontentslabel]{1.25cm}\color{ocre}} % Chapter number
	{}  
	{\color{ocre}\normalsize\sffamily\bfseries\;\titlerule*[.5pc]{.}\;\thecontentspage} % Page number

% Section text styling
\titlecontents{lsection}
	[0em] % Left indentation
	{\sffamily\small} % Spacing and font options for sections
	{\contentslabel[\thecontentslabel]{1.25cm}} % Section number
	{}
	{}

% Subsection text styling (note these aren't shown by default, display them by searchings this file for tocdepth and reading the commented text)
\titlecontents{lsubsection}
	[.5em] % Left indentation
	{\sffamily\footnotesize} % Spacing and font options for subsections
	{\contentslabel[\thecontentslabel]{1.25cm}}
	{}
	{}

%----------------------------------------------------------------------------------------
%	HEADERS AND FOOTERS
%----------------------------------------------------------------------------------------


\usepackage{fancyhdr} % Required for header and footer configuration

\pagestyle{fancy}
\renewcommand{\chaptermark}[1]{\markboth{\sffamily\normalsize\bfseries\ \thechapter.\ #1}{}} % Chapter text font settings
\renewcommand{\sectionmark}[1]{\markright{\sffamily\normalsize\thesection\hspace{5pt}#1}{}} % Section text font settings
\fancyhf{} \fancyhead[LE,RO]{\sffamily\normalsize\thepage} % Font setting for the page number in the header
\fancyhead[LO]{\rightmark} % Print the nearest section name on the left side of odd pages
\fancyhead[RE]{\leftmark} % Print the current chapter name on the right side of even pages

\fancyfoot[L]{Jason LAPEYRONNIE}
\fancyfoot[R]{\href{http://mathoutils.fr}{http://mathoutils.fr}} % Uncomment to include a footer

\renewcommand{\headrulewidth}{0.5pt} % Thickness of the rule under the header
\renewcommand{\footrulewidth}{0.5pt} % Thickness of the rule under the header

\fancypagestyle{plain}{% Style for when a plain pagestyle is specified
	\fancyhead{}\renewcommand{\headrulewidth}{0pt}%
}

% Removes the header from odd empty pages at the end of chapters
\makeatletter
\renewcommand{\cleardoublepage}{
\clearpage\ifodd\c@page\else
\hbox{}
\vspace*{\fill}
\thispagestyle{empty}
\newpage
\fi}

%----------------------------------------------------------------------------------------
%	THEOREM STYLES
%----------------------------------------------------------------------------------------

\usepackage{amsmath,amsfonts,amssymb,amsthm} % For math equations, theorems, symbols, etc

\newcommand{\intoo}[2]{\mathopen{]}#1\,;#2\mathclose{[}}
\newcommand{\ud}{\mathop{\mathrm{{}d}}\mathopen{}}
\newcommand{\intff}[2]{\mathopen{[}#1\,;#2\mathclose{]}}
\renewcommand{\qedsymbol}{$\blacksquare$}
\newtheorem{notation}{Notation}[section]

% Boxed/framed environments
\newtheoremstyle{ocrenumbox}% Theorem style name
{0pt}% Space above
{0pt}% Space below
{\normalfont}% Body font
{}% Indent amount
{\small\bf\sffamily\color{ocre}}% Theorem head font
{\;:\;}% Punctuation after theorem head
{0.25em}% Space after theorem head
{\small\sffamily\color{ocre}\thmname{#1}\nobreakspace\thmnumber{\@ifnotempty{#1}{}\@upn{#2}}% Theorem text (e.g. Theorem 2.1)
\thmnote{\nobreakspace\the\thm@notefont\sffamily\bfseries\color{black}---\nobreakspace#3}} % Optional theorem note

\newtheoremstyle{blacknumex}% Theorem style name
{5pt}% Space above
{10pt}% Space below
{\normalfont}% Body font
{} % Indent amount
{\small\bf\sffamily}% Theorem head font
{\;:\;}% Punctuation after theorem head
{0.25em}% Space after theorem head
{\small\sffamily{\tiny\ensuremath{\blacksquare}}\nobreakspace\thmname{#1}\nobreakspace\thmnumber{\@ifnotempty{#1}{}\@upn{#2}}% Theorem text (e.g. Theorem 2.1)
\thmnote{\nobreakspace\the\thm@notefont\sffamily\bfseries---\nobreakspace#3}}% Optional theorem note

\newtheoremstyle{blacknumexo}% Theorem style name
{15pt}% Space above
{10pt}% Space below
{\normalfont}% Body font
{} % Indent amount
{\small\bf\sffamily}% Theorem head font
{}% Punctuation after theorem head
{0.5em}% Space after theorem head
{\small\sffamily{\ensuremath{\blacktriangleright}}\nobreakspace\thmname{#1}\nobreakspace\thmnumber{\@ifnotempty{#1}{}\@upn{#2}}% Theorem text (e.g. Theorem 2.1)
\thmnote{\nobreakspace\the\thm@notefont\sffamily\bfseries---\nobreakspace#3} \\}% Optional theorem note



\newtheoremstyle{blacknumbox} % Theorem style name
{0pt}% Space above
{5pt}% Space below
{}% Body font
{}% Indent amount
{\large\bf\sffamily}% Theorem head font
{\;:\;}% Punctuation after theorem head
{0.25em}% Space after theorem head
{\small\sffamily\thmname{#1}\nobreakspace\thmnumber{\@ifnotempty{#1}{}\@upn{#2}}% Theorem text (e.g. Theorem 2.1)
\thmnote{\nobreakspace\the\thm@notefont\sffamily\bfseries---\nobreakspace#3}}% Optional theorem note

% Non-boxed/non-framed environments
\newtheoremstyle{ocrenum}% Theorem style name
{5pt}% Space above
{5pt}% Space below
{\normalfont}% Body font
{}% Indent amount
{\small\bf\sffamily\color{ocre}}% Theorem head font
{\;:\;}% Punctuation after theorem head
{0.25em}% Space after theorem head
{\small\sffamily\color{ocre}\thmname{#1}\nobreakspace\thmnumber{\@ifnotempty{#1}{}\@upn{#2}}% Theorem text (e.g. Theorem 2.1)
\thmnote{\nobreakspace\the\thm@notefont\sffamily\bfseries\color{black}---\nobreakspace#3}} % Optional theorem note
\makeatother

% Defines the theorem text style for each type of theorem to one of the three styles above
\newcounter{dummy} 
\newcounter{thm}
\newcounter{correction}
\newcounter{qst}
\theoremstyle{ocrenumbox}
\newtheorem{theoremeT}[dummy]{Théorème}
\newtheorem{exerciseT}{Propriété}
\newtheorem{principeT}{Principe}
\theoremstyle{blacknumex}
\newtheorem{exampleT}{Exemple}
\theoremstyle{blacknumexo}
\newtheorem{exo}[thm]{Exercice}
\newtheorem{corr}[correction]{Correction}
\newtheorem{quest}[qst]{Question}
\theoremstyle{blacknumbox}
\newtheorem{vocabulary}{Vocabulary}[section]
\newtheorem{definitionT}{Définition}
\newtheorem{corollaryT}[dummy]{Corollary}
\theoremstyle{ocrenum}
\newtheorem{proofT}[dummy]{Démonstration}


%----------------------------------------------------------------------------------------
%	DEFINITION OF COLORED BOXES
%----------------------------------------------------------------------------------------

\RequirePackage[framemethod=default]{mdframed} % Required for creating the theorem, definition, exercise and corollary boxes

% Theorem box
\newmdenv[skipabove=7pt,
skipbelow=7pt,
backgroundcolor=black!5,
linecolor=ocre,
innerleftmargin=5pt,
innerrightmargin=5pt,
innertopmargin=10pt,
leftmargin=0cm,
rightmargin=0cm,
innerbottommargin=5pt]{tBox}

%Proposition box	  
\newmdenv[skipabove=7pt,
skipbelow=7pt,
rightline=false,
leftline=true,
topline=false,
bottomline=false,
backgroundcolor=ocre!10,
linecolor=ocre,
innerleftmargin=5pt,
innerrightmargin=5pt,
innertopmargin=10pt,
innerbottommargin=3pt,
leftmargin=0cm,
rightmargin=0cm,
linewidth=4pt]{eBox}	

% Definition box
\newmdenv[skipabove=7pt,
backgroundcolor=ocre!4,
skipbelow=7pt,
rightline=false,
leftline=true,
topline=false,
bottomline=false,
linecolor=ocre,
innerleftmargin=5pt,
innerrightmargin=5pt,
innertopmargin=10pt,
leftmargin=0cm,
rightmargin=0cm,
linewidth=4pt,
innerbottommargin=5pt]{dBox}	

% Corollary box
\newmdenv[skipabove=7pt,
skipbelow=7pt,
rightline=false,
leftline=true,
topline=false,
bottomline=false,
linecolor=gray,
backgroundcolor=black!5,
innerleftmargin=5pt,
innerrightmargin=5pt,
innertopmargin=5pt,
leftmargin=0cm,
rightmargin=0cm,
linewidth=4pt,
innerbottommargin=5pt]{cBox}

\newmdenv[skipabove=7pt,
skipbelow=7pt,
backgroundcolor=black!5,
innerleftmargin=5pt,
topline=false,
bottomline=false,
rightline=false,
leftline=false,
innerrightmargin=5pt,
innertopmargin=5pt,
leftmargin=0cm,
rightmargin=0cm,
innerbottommargin=5pt]{xBox}

% Creates an environment for each type of theorem and assigns it a theorem text style from the "Theorem Styles" section above and a colored box from above
\newenvironment{theorem}{\begin{tBox}\begin{theoremeT}}{\end{theoremeT}\end{tBox}}

\newenvironment{exo2}{\noindent \begin{exo}\item\relax \noindent \begin{eBox}\item\relax}{\end{eBox}\end{exo}}


\newenvironment{proposition}{\begin{eBox}\begin{exerciseT}}{\hfill{\color{ocre}}\end{exerciseT}\end{eBox}}		

\newenvironment{principe}{\begin{eBox}\begin{principeT}}{\hfill{\color{ocre}}\end{principeT}\end{eBox}}	
		  
\newenvironment{definition}{\begin{dBox}\begin{definitionT}}{\end{definitionT}\end{dBox}}	

\newenvironment{example}{\begin{xBox}\begin{exampleT}}{\hfill{\tiny\ensuremath{\blacksquare}}\end{exampleT}\end{xBox}}

\newenvironment{demonstration}{\begin{proofT}}{\hfill{\tiny\ensuremath{\square}}\end{proofT}}		
\newenvironment{corollary}{\begin{cBox}\begin{corollaryT}}{\end{corollaryT}\end{cBox}}	

%----------------------------------------------------------------------------------------
%	REMARK ENVIRONMENT
%----------------------------------------------------------------------------------------

\newenvironment{remark}{\par\vspace{5pt}\small % Vertical white space above the remark and smaller font size
\begin{list}{}{
\leftmargin=25pt % Indentation on the left
\rightmargin=15pt}\item\ignorespaces % Indentation on the right
\makebox[-2.5pt]{\begin{tikzpicture}[overlay]
\node[draw=ocre!60,line width=1pt,circle,fill=ocre!25,font=\sffamily\bfseries,inner sep=2pt,outer sep=0pt] at (-15pt,0pt){\textcolor{ocre}{R}};\end{tikzpicture}} % Orange R in a circle
\advance\baselineskip -1pt}{\end{list}\vskip5pt} % Tighter line spacing and white space after remark

%----------------------------------------------------------------------------------------
%	SECTION NUMBERING IN THE MARGIN
%----------------------------------------------------------------------------------------

\makeatletter
\renewcommand{\@seccntformat}[1]{\llap{\textcolor{ocre}{\csname the#1\endcsname}\hspace{1em}}}                    
\renewcommand{\section}{\@startsection{section}{1}{\z@}
{-4ex \@plus -1ex \@minus -.4ex}
{1ex \@plus.2ex }
{\normalfont\LARGE\sffamily\bfseries}}
\renewcommand{\subsection}{\@startsection {subsection}{2}{\z@}
{-3ex \@plus -0.1ex \@minus -.4ex}
{0.5ex \@plus.2ex }
{\normalfont\sffamily\bfseries}}
\renewcommand{\subsubsection}{\@startsection {subsubsection}{3}{\z@}
{-2ex \@plus -0.1ex \@minus -.2ex}
{.2ex \@plus.2ex }
{\normalfont\small\sffamily\bfseries}}                        
\renewcommand\paragraph{\@startsection{paragraph}{4}{\z@}
{-2ex \@plus-.2ex \@minus .2ex}
{.1ex}
{\normalfont\small\sffamily\bfseries}}

%----------------------------------------------------------------------------------------
%	PART HEADINGS
%----------------------------------------------------------------------------------------

% Numbered part in the table of contents
\newcommand{\@mypartnumtocformat}[2]{%
	\setlength\fboxsep{0pt}%
	\noindent\colorbox{ocre!20}{\strut\parbox[c][.7cm]{\ecart}{\color{ocre!70}\Large\sffamily\bfseries\centering#1}}\hskip\esp\colorbox{ocre!40}{\strut\parbox[c][.7cm]{\linewidth-\ecart-\esp}{\Large\sffamily\centering#2}}%
}

% Unnumbered part in the table of contents
\newcommand{\@myparttocformat}[1]{%
	\setlength\fboxsep{0pt}%
	\noindent\colorbox{ocre!40}{\strut\parbox[c][.7cm]{\linewidth}{\Large\sffamily\centering#1}}%
}

\newlength\esp
\setlength\esp{4pt}
\newlength\ecart
\setlength\ecart{1.2cm-\esp}
\newcommand{\thepartimage}{}%
\newcommand{\partimage}[1]{\renewcommand{\thepartimage}{#1}}%
\def\@part[#1]#2{%
\ifnum \c@secnumdepth >-2\relax%
\refstepcounter{part}%
\addcontentsline{toc}{part}{\texorpdfstring{\protect\@mypartnumtocformat{\thepart}{#1}}{\partname~\thepart\ ---\ #1}}
\else%
\addcontentsline{toc}{part}{\texorpdfstring{\protect\@myparttocformat{#1}}{#1}}%
\fi%
\startcontents%
\markboth{}{}%
{\thispagestyle{empty}%
\begin{tikzpicture}[remember picture,overlay]%
\node at (current page.north west){\begin{tikzpicture}[remember picture,overlay]%	
\fill[ocre!20](0cm,0cm) rectangle (\paperwidth,-\paperheight);
\node[anchor=north] at (4cm,-3.25cm){\color{ocre!40}\fontsize{220}{100}\sffamily\bfseries\thepart}; 
\node[anchor=south east] at (\paperwidth-1cm,-\paperheight+1cm){\parbox[t][][t]{8.5cm}{
\printcontents{l}{0}{\setcounter{tocdepth}{1}}% The depth to which the Part mini table of contents displays headings; 0 for chapters only, 1 for chapters and sections and 2 for chapters, sections and subsections
}};
\node[anchor=north east] at (\paperwidth-1.5cm,-3.25cm){\parbox[t][][t]{15cm}{\strut\raggedleft\color{white}\fontsize{30}{30}\sffamily\bfseries#2}};
\end{tikzpicture}};
\end{tikzpicture}}%
\@endpart}
\def\@spart#1{%
\startcontents%
\phantomsection
{\thispagestyle{empty}%
\begin{tikzpicture}[remember picture,overlay]%
\node at (current page.north west){\begin{tikzpicture}[remember picture,overlay]%	
\fill[ocre!20](0cm,0cm) rectangle (\paperwidth,-\paperheight);
\node[anchor=north east] at (\paperwidth-1.5cm,-3.25cm){\parbox[t][][t]{15cm}{\strut\raggedleft\color{white}\fontsize{30}{30}\sffamily\bfseries#1}};
\end{tikzpicture}};
\end{tikzpicture}}
\addcontentsline{toc}{part}{\texorpdfstring{%
\setlength\fboxsep{0pt}%
\noindent\protect\colorbox{ocre!40}{\strut\protect\parbox[c][.7cm]{\linewidth}{\Large\sffamily\protect\centering #1\quad\mbox{}}}}{#1}}%
\@endpart}
\def\@endpart{\vfil\newpage
\if@twoside
\if@openright
\null
\thispagestyle{empty}%
\newpage
\fi
\fi
\if@tempswa
\twocolumn
\fi}

%----------------------------------------------------------------------------------------
%	CHAPTER HEADINGS
%----------------------------------------------------------------------------------------

% A switch to conditionally include a picture, implemented by Christian Hupfer
\newif\ifusechapterimage
\usechapterimagetrue
\newcommand{\thechapterimage}{}%
\newcommand{\chapterimage}[1]{\ifusechapterimage\renewcommand{\thechapterimage}{#1}\fi}%
\newcommand{\autodot}{.}
\def\@makechapterhead#1{%
{\parindent \z@ \raggedright \normalfont
\ifnum \c@secnumdepth >\m@ne
\if@mainmatter
\begin{tikzpicture}[remember picture,overlay]
\node at (current page.north west)
{\begin{tikzpicture}[remember picture,overlay]
\node[anchor=north west,inner sep=0pt] at (0,0) {\ifusechapterimage\includegraphics[width=\paperwidth]{\thechapterimage}\fi};
\draw[anchor=west] (\Gm@lmargin,-3cm) node [line width=2pt,rounded corners=15pt,draw=ocre,fill=white,fill opacity=0.5,inner sep=15pt]{\strut\makebox[22cm]{}};
\draw[anchor=west] (\Gm@lmargin+.3cm,-3cm) node {\huge\sffamily\bfseries\color{black}\thechapter\autodot~#1\strut};
\end{tikzpicture}};
\end{tikzpicture}
\else
\begin{tikzpicture}[remember picture,overlay]
\node at (current page.north west)
{\begin{tikzpicture}[remember picture,overlay]
\node[anchor=north west,inner sep=0pt] at (0,0) {\ifusechapterimage\includegraphics[width=\paperwidth]{\thechapterimage}\fi};
\draw[anchor=west] (\Gm@lmargin,-3cm) node [line width=2pt,rounded corners=15pt,draw=ocre,fill=white,fill opacity=0.5,inner sep=15pt]{\strut\makebox[22cm]{}};
\draw[anchor=west] (\Gm@lmargin+.3cm,-3cm) node {\huge\sffamily\bfseries\color{black}#1\strut};
\end{tikzpicture}};
\end{tikzpicture}
\fi\fi\par\vspace*{50\p@}}}

%-------------------------------------------

\def\@makeschapterhead#1{%
\begin{tikzpicture}[remember picture,overlay]
\node at (current page.north west)
{\begin{tikzpicture}[remember picture,overlay]
\node[anchor=north west,inner sep=0pt] at (0,0) {\ifusechapterimage\includegraphics[width=\paperwidth]{\thechapterimage}\fi};
\draw[anchor=west] (\Gm@lmargin,-3cm) node [line width=2pt,rounded corners=15pt,draw=ocre,fill=white,fill opacity=0.5,inner sep=15pt]{\strut\makebox[22cm]{}};
\draw[anchor=west] (\Gm@lmargin+.3cm,-3cm) node {\huge\sffamily\bfseries\color{black}#1\strut};
\end{tikzpicture}};
\end{tikzpicture}
\par\vspace*{50\p@}}
\makeatother

%----------------------------------------------------------------------------------------
%	LINKS
%----------------------------------------------------------------------------------------

\usepackage{hyperref}
\hypersetup{hidelinks,backref=true,pagebackref=true,hyperindex=true,colorlinks=false,breaklinks=true,urlcolor=ocre,bookmarks=true,bookmarksopen=false}

\usepackage{bookmark}
\bookmarksetup{
open,
numbered,
addtohook={%
\ifnum\bookmarkget{level}=0 % chapter
\bookmarksetup{bold}%
\fi
\ifnum\bookmarkget{level}=-1 % part
\bookmarksetup{color=ocre,bold}%
\fi
}
}

\renewcommand*\thesection{\arabic{section}}

\newcommand*{\coord}[3]{% 
  \ensuremath{\overrightarrow{#1}\, 
    \begin{pmatrix} 
      #2\\ 
      #3 
    \end{pmatrix}}}
    
  \newcommand*{\coordb}[2]{% 
  \ensuremath{ 
    \begin{pmatrix} 
      #1\\ 
      #2 
    \end{pmatrix}}}

\newcommand*{\coorde}[4]{% 
  \renewcommand{\arraystretch}{1}\ensuremath{\overrightarrow{#1}\, 
    \begin{pmatrix} 
      #2\\ 
      #3 \\
      #4
    \end{pmatrix}}}    
  \newcommand*{\coordbe}[3]{% 
 \renewcommand{\arraystretch}{1} \ensuremath{ 
    \begin{pmatrix} 
      #1\\ 
      #2 \\
      #3
    \end{pmatrix}}}  
    
\newcommand{\Card}{\mathrm{Card}}


\DeclareExerciseCollection[topic=cpx01]{cpx01}
\DeclareExerciseCollection[topic=cpx02]{cpx02}
\DeclareExerciseCollection[topic=cpx03]{cpx03}
\DeclareExerciseCollection[topic=cpx04]{cpx04}


\begin{document}

\chapterimage{../../Pictures/background}
\chapter{Nombres complexes : forme algébrique}


\section{Ensemble $\mathbb{C}$ des nombres complexes}

\begin{definition}Il existe un ensemble de nombres appelé \textbf{nombres complexes} et noté $\mathbb{C}$ tel que
\begin{itemize}
\item L'ensemble des réels $\mathbb{R}$ est inclus dans $\mathbb{C}$.
\item Il existe un nombre complexe, noté $i$, et tel que $i^2=-1$.
\item L'addition et la multiplication des réels se prolonge "naturellement" dans l'ensemble des complexes.
\item Pour tout nombre complexe $z$, il existe un unique couple de réels $(a,b)$ tel que $z=a+ib$.
\begin{itemize}
\item $a$ est appelée \textbf{partie réelle} de $z$, notée $Re(z)$.
\item $b$ est appelée \textbf{partie imaginaire} de $z$, notée $Im(z)$.
\item L'écriture $z=a+ib$ est la \textbf{forme algébrique} de $z$.
\end{itemize}
\end{itemize}\end{definition}

\begin{example}$2+3i$, $7i$, $\sqrt{5}$, $\frac{1}{2}-\frac{4}{5}i$ sont des nombres complexes.\end{example}

\begin{definition}Soit $z$ un nombre complexe.
\begin{itemize}
\item $z$ est un nombre réel si et seulement si sa partie imaginaire est nulle.
\item Si $Re(z)=0$, on dit que $z$ est un nombre imaginaire pur. On note $z \in i\mathbb{R}$.
\end{itemize}\end{definition}

\begin{example}$2i$ est un nombre imaginaire pur.\end{example}

L'addition et la multiplication des nombres réels se prolongent naturellement pour les nombres complexes. Le nombre $i$ joue le rôle de facteur. Il ne faut toutefois pas oublier que $i^2=-1$, notamment pour la multiplication. Ainsi, on peut développer, factoriser, et l'on retrouve évidemment les identités remarquables.

\begin{example}Soit $z=2+5i$ et $z'=1-3i$. Alors $z +z'=2+5i+1-3i=3+2i$ et
\[zz' = (2+5i)(1-3i)=2 \times 1 + 5i \times 1 + 2 \times (-3i) + 5i \times (-3i)=2+5i-6i -15i^2=17-i.\]
Enfin, $z^2=(2+5i)^2=4^2+2\times 2 \times 5i + (5i)^2 = 16+20i-25=-9+20i$.\end{example}

\begin{proposition}Soit $z$ un nombre complexe. On a $z=0$ si et seulement si $Re(z)=0$ et $Im(z)=0$.\end{proposition}

\begin{demonstration}
D'une part, si $Re(z)=0$ et $Im(z)=0$, alors $z = Re(z) + i \times Im(z) = 0 + i0 = 0$.

D'autre part, si $z=0$, notons $a$ et $b$ les parties réelles et imaginaires de $z$. Ainsi, $a+ib=0$. 

Nous allons raisonner par l'absurde pour montrer que nécessairement, $a$ et $b$ valent tous deux 0. Pour cela, nous allons supposer le contraire et aboutir à une absurdité, ce qui nous permettra de conclure que l'hypothèse faite au départ était fausse. Supposons donc que soit $a$, soit $b$ ne soit pas nul.
\begin{itemize}
\item Si $b=0$, puisque $a+ib=0$, on obtient également que $a=0$, ce qui est contraire à ce que nous avons supposé. Ainsi, on a forcément $b \neq 0$.
\item Si $b\neq 0$, alors, puisque $a+ib=0$, on a donc $i=-\dfrac{a}{b}$. $i$ serait donc réel. C'est absurde.\end{itemize}

Ainsi, $a=b=0$.\end{demonstration}

\begin{proposition}Soit $z$ et $z'$ deux complexes. Alors $z = z'$ si et seulement si $Re(z)=Re(z')$  et $Im(z)=Im(z')$.\end{proposition}

\begin{demonstration}Notons $z=a+ib$ et $z'=a'+ib'$. On a $z=z'$ si et seulement si $a+ib=a'+ib'$ si et seulement si $a-a'+i(b-b')=0$. D'après le premier point, c'est équivalent à dire que $a-a'=0$ et $b-b'=0$. 

Finalement, $z=z'$ si et seulement si $a=a'$ et $b=b'$.\end{demonstration}

Cette propriété nous permet notamment de résoudre des équations dans $\mathbb{C}$.

\begin{example}Résoudre l'équation $2z+i-3=5i+4z-7$, d'inconnue $z \in \mathbb{C}$.

Soit $z \in \mathbb{C}$. On a
\[2z+i-3=5i+4z-7  \Leftrightarrow  2z-4z = 5i - 7 +3 - i  \Leftrightarrow -2z = -4 +4i  \Leftrightarrow  z = 2 - 2i\]

Ainsi, $S=\{2-2i\}$.

\end{example}


\section{Conjugué d'un nombre complexe}

\subsection{Définition et propriétés}

\begin{definition}Soit $z = a+ib$ un nombre complexe.

Le \textbf{conjugué} de $z$ est le nombre complexe noté $\overline{z}$ et qui vaut $\overline{z}=a-ib$.\end{definition}

\begin{example}Le conjuqué de $7+3i$ est $7-3i$. Le conjugué de $2i$ est $-2i$.\end{example}


\begin{proposition}Soit $z \in \mathbb{C}$.
\begin{itemize}
\item $z$ est réel si et seulement si $\overline{z}=z$ ;
\item $z$ est imaginaire pur si et seulement si $\overline{z}=-z$.
\end{itemize}\end{proposition}

\begin{demonstration}[Premier point]Soit $z=a+ib$ avec $a$ et $b$ des réels. On a donc $\overline{z}=a-ib$.
\begin{itemize}
\item Si $z=\overline{z}$, alors $a+ib=a-ib$ ce qui implique que $2ib=0$ et donc $b=0$. $z$ est donc réel.
\item Si $z$ est réel, alors $b=0$. Ainsi, $z=a$ et $\overline{z}=a$. On a bien $z=\overline{z}$.
\end{itemize}\end{demonstration}

La démonstration du deuxième point fait l'objet d'un exercice.

\begin{proposition}Soit $z$ et $z'$ deux nombres complexes. On a les propriétés suivantes
\[\overline{\overline{z}} = z \qquad \qquad \overline{z+z'}=\overline{z}+\overline{z'} \qquad \qquad \overline{zz'}=\overline{z}\times \overline{z'}\]
\end{proposition}

\begin{demonstration}Notons $z=a+ib$ et $z'=a'+ib'$.

On a $\overline{\overline{z}} = \overline{a-ib}=a-(-ib)=a+ib=z$.

Par ailleurs, $\overline{z+z'} = \overline{a+ib+a'+ib'}=\overline{a+a'+i(b+b')}=a+a'-i(b+b')=a-ib+a'-ib'=z-z'$.

Pour montrer la propriété sur le produit, on raisonne en deux temps.
\begin{itemize}
\item D'une part, \[\overline{z} \times \overline{z'}=(a-ib)(a'-ib')=aa'-iab'-ia'b+i^2bb'=aa'-bb'-i(ab'+a'b).\]
 \item D'autre part, 
 \[zz'=(a+ib)(a'+ib')=aa'+iab'+ia'b+i^2bb'=aa'-bb'+i(ab'+a'b).\]
\end{itemize}
 On a bien $\overline{zz'}=\overline{z}\times \overline{z'}$.\end{demonstration}



\begin{proposition}Soit $z$ un nombre complexe et $n$ un entier naturel non nul. Alors $\overline{z^n}=\overline{z}^n$.\end{proposition}

\begin{demonstration}Soit $z$ un nombre complexe. Nous allons démontrer cette propriété par récurrence.

Pour tout entier naturel non nul $n$, on considère la proposition $\mathcal{P}(n)$ : « $\overline{z^n}=\overline{z}^n$ ».

\begin{itemize}
\item \textbf{Initialisation} : Pour $n=1$, on a $\overline{z^1}=\overline{z}=\overline{z}^1$. $\mathcal{P}(1)$ est vraie.
\item \textbf{Hérédité} : Soit $n$ un entier naturel non nul. Supposons que $\mathcal{P}(n)$ est vraie. On a donc $\overline{z^n}=\overline{z}^n$.

Or, $\overline{z^{n+1}}=\overline{z^n  \times z}$. La propriété précédente nous assure donc que $\overline{z^{n+1}}=\overline{z^n} \times \overline{z}$. 

Or, par hypothèse de récurrence, $\overline{z^n}=\overline{z}^n$.

Ainsi, $\overline{z^{n+1}}=\overline{z}^n \times \overline{z}=\overline{z}^{n+1}$. $\mathcal{P}(n+1)$ est donc vraie.
\item \textbf{Conclusion} : $\mathcal{P}(1)$ est vraie et $\mathcal{P}$ est héréditaire. Par récurrence, $\mathcal{P}(n)$ est vraie pour tout entier naturel non nul $n$.
\end{itemize}\end{demonstration}

\subsection{Résolution d'équation en $z$ et $\overline{z}$}

Lorsque des équations font intervenir $z$ et $\overline{z}$, il est préférable d'écrire ces nombres sous forme algébrique, puis d'identifier les parties réelles et imaginaires de chaque membre de l'équation.

\begin{example}On souhaite résoudre l'équation $(E) \,: \,2z + 3i -5 = 3i\overline{z}-2i+5$, d'inconnue $z \in \mathbb{C}$.

On pose alors $z=a+ib$. On a donc $\overline{z}=a-ib$. Ainsi,

\begin{eqnarray*}
2z + 3i -5 = 3\overline{z}-2i+7 & \Leftrightarrow & 2(a+ib)+3i-5=3i(a-ib)-2i+5 \\
& \Leftrightarrow & 2a+2ib+3i-5=3ia+3b-2i+5=0 \\
& \Leftrightarrow & 2a-3b-5-5 +i(2b+3-3a+2)=0 \\
& \Leftrightarrow & 2a -3b-10 + i (2b-3a+5)=0
\end{eqnarray*}
Or, un complexe est nul si et seulement si sa partie réelle et sa partie imaginaires sont nulles. On obtient donc le système suivant.
\[(E) \quad \Leftrightarrow \quad \left\{\begin{array}{rcl}
2a -3b-10 &=& 0\\
-3a+2b+5 &=& 0
\end{array}\right.\]

On peut alors résoudre ce système avec la méthode de son choix. Par exemple, en procédant par combinaison. On multiplie la première ligne par 3 et la deuxième par 2.
\[(E) \quad \Leftrightarrow \quad\left\{\begin{array}{rcl}
6a -9b-30 &=& 0\\
-6a+4b+10 &=& 0
\end{array}\right.\]

On remplace $L_2$ par $L_2+L_1$
\[(E) \quad \Leftrightarrow \quad\left\{\begin{array}{rcl}
6a -9b-30 &=& 0\\
-5b-20 &=& 0
\end{array}\right.\]

On obtient donc
\[(E) \quad \Leftrightarrow \quad\left\{\begin{array}{rcl}
6a -9\times (-4)-30 &=& 0\\
b &=& -4
\end{array}\right. \Leftrightarrow \quad\left\{\begin{array}{rcl}
6a +6 &=& 0\\
b &=& -4
\end{array}\right.  \Leftrightarrow \quad\left\{\begin{array}{rcl}
a &=& -1\\
b &=& -4
\end{array}\right.\]
On a donc $S=\{ -1-4i\}$.
\end{example}

\section{Division dans $\mathbb{C}$}

\begin{definition}Soit $z$ un nombre complexe non nul. Il existe un unique complexe $z'$ tel que $zz'=1$. 

$z'$ est appelé inverse de $z$ et est noté $\dfrac{1}{z}$.\end{definition}

\begin{example} On a $(1+2i)\left(\dfrac{1}{5}-\dfrac{2}{5}i\right)=\dfrac{1}{5}+\dfrac{2}{5}i-\dfrac{2}{5}i+\dfrac{4}{5}=1$. On a donc $\dfrac{1}{5}-\dfrac{2}{5}i=\dfrac{1}{1+2i}$.\end{example}


\begin{definition}Soit $z$ un complexe et $\overline{z}$ un complexe non nul. 

On définit le quotient de $z$ par $z'$ par  $\dfrac{z}{z'}=z \times \dfrac{1}{z'}$.\end{definition}

\begin{example} Calculons $\dfrac{3+i}{1+2i}$. On a alors $\dfrac{3+i}{1+2i} = (3+i) \times \dfrac{1}{1+2i}$.

Or, d'après l'exemple précédent, $\dfrac{1}{1+2i}=\dfrac{1}{5}-\dfrac{2}{5}i$. Ainsi,
\[ \dfrac{3+i}{1+2i} = (3+i)\left(\dfrac{1}{5}-\dfrac{2}{5}i\right) = \dfrac{3}{5}-\dfrac{6}{5}i+\dfrac{i}{5}-\dfrac{2}{5}i^2=1-i.\]\end{example}


Dans un précédent exercice sur les conjugués, nous avons pu exprimer le produit $z\overline{z}$ et constater que ce produit était en fait un réel.

Pour mettre sous forme algébrique un quotient, on multiplie le numérateur et le dénominateur par le conjugué du dénominateur. Ainsi, le dénominateur du quotient obtenu sera un nombre réel.

\begin{example}On souhaite exprimer $\dfrac{4+3i}{3-i}$ sous forme algébrique. 

Le dénominateur est $3-i$, dont le conjugué est $3+i$. 

On multiplie donc numérateur et dénominateur par $3+i$.
\[ \dfrac{4+3i}{3-i}= \dfrac{(4+3i)(3+i)}{(3-i)(3+i)}=\dfrac{12+4i+9i+3i^2}{3^2-i^2}=\dfrac{12+13i-3}{9-(-1)}=\dfrac{9}{10}+\dfrac{13}{10}i.\]\vspace{-0.5cm}\end{example}

\begin{example}On souhaite résoudre l'équation $2z+3-i=(1+i)z-2+i$. Soit donc $z \in \mathbb{C}$
\begin{eqnarray*}
2z+3-i=(1+i)z-2+i & \Leftrightarrow & 2z-(1+i)z=-2+i-3+i \\
& \Leftrightarrow & (1-i)z = -5+2i \\
& \Leftrightarrow & z = \dfrac{-5+2i}{1-i}
\end{eqnarray*}

Il reste alors à écrire $\dfrac{-5+2i}{1-i}$ sous forme algébrique.
\[\dfrac{-5+2i}{1-i}=\dfrac{(-5+2i)(1+i)}{(1-i)(1+i)}=\dfrac{-5-5i+2i-2}{1^2-i^2}=\dfrac{-7-3i}{2}=-\dfrac{7}{2}-\dfrac{3}{2}i\]
L'unique solution de l'équation $2z+3-i=(1+i)z-2+i$ est donc $-\dfrac{7}{2}-\dfrac{3}{2}i$.\end{example}

\begin{proposition}Soit $z$ un complexe et $z'$ un complexe non nul. Alors
\[ \overline{\left(\dfrac{1}{z'}\right)} = \dfrac{1}{\overline{z'}} \quad \text{et} \overline{\left(\dfrac{z}{z'}\right)} = \dfrac{\overline{z}}{\overline{z'}}.\]\end{proposition}

\begin{demonstration}Soit $z$ un complexe non nul. 

D'après la propriété sur le produit des conjugués, on a $\overline{z} \times \overline{\dfrac{1}{z}}=\overline{z \times \dfrac{1}{z}} = \overline{1}=1$ et donc
$\overline{\left(\dfrac{1}{z}\right)} = \dfrac{1}{\overline{z}}$.

Par ailleurs, $\overline{\left(\dfrac{z}{z'}\right)} = \overline{z \times \dfrac{1}{z'}} = \overline{z} \times \overline{\left(\dfrac{1}{z'}\right)}$
d'après la proposition sur le conjugué du produit. 

Mais alors, d'après la propriété précédemment démontrée, $\overline{\left(\dfrac{z}{z'}\right)} = \overline{z} \times \dfrac{1}{\overline{z'}}= \dfrac{\overline{z}}{\overline{z'}}$.
\end{demonstration}

\section{Équations du second degré à coefficients réels}

\begin{proposition}Soit $a$ un réel strictement positif.

Les solutions de l'équation $z^2=-a$ d'inconnue complexe $z$ sont les réels $i\sqrt{a}$ et $-i\sqrt{a}$.\end{proposition}

\begin{demonstration}On a en effet $z^2+a=z^2-i^2\times a = z^2 - (i\sqrt{a})^2=(z-i\sqrt{a})(z+i\sqrt{a})$.

Ainsi, $z^2+a=0$ si et seulement si $z-i\sqrt{a}=0$ ou $z+i\sqrt{a}=0$, d'où le résultat.\end{demonstration}


\begin{example}Les solutions de l'équation $z^2=-4$ sont $2i$ et $-2i$.\end{example}

L'exemple que nous venons de voir est un peu simpliste mais permet de vous rafraîchir la mémoire. Passons donc aux choses sérieuses.

\newpage

\begin{theorem}Soit $a$, $b$ et $c$ des \textbf{réels}. Soit $P : z \mapsto az^2+bz+c$.

Notons $\Delta = b^2-4ac$ le discriminant du polynôme $P$.
\begin{itemize}
\item Si $\Delta >0$, alors $P$ admet deux racines réelles distinctes : $\dfrac{-b-\sqrt{\Delta}}{2a}$ et $\dfrac{-b+\sqrt{\Delta}}{2a}$ ;
\item Si $\Delta = 0$, alors $P$ admet une racine réelle double : $-\dfrac{b}{2a}$ ;
\item Si $\Delta <0$, alors $P$ admet deux racines complexes conjuguées : $\dfrac{-b-i\sqrt{-\Delta}}{2a}$ et $\dfrac{-b+i\sqrt{-\Delta}}{2a}$.
\end{itemize}
\end{theorem}

\begin{demonstration}
Écrivons $P$ sous forme canonique. Pour tout complexe $z$,
\[az^2+bz+c = a\left(z+\dfrac{b}{2a}\right)^2- \dfrac{\Delta}{4a} = a\left(\left(z+\dfrac{b}{2a}\right)^2- \dfrac{\Delta}{4a^2}\right).\]

\begin{itemize}
\item Si $\Delta > 0$, alors, pour tout complexe $z$, 
\[az^2+bz+c = a\left(\left(z+\dfrac{b}{2a}\right)^2- \left(\dfrac{\sqrt{\Delta}}{2a}\right)^2\right) = a\left(z+\dfrac{b-\sqrt{\Delta}}{2a}\right)\left(z+\dfrac{b+\sqrt{\Delta}}{2a}\right).\]
Ainsi, $P$ possède deux racines réelles qui sont $\dfrac{-b-\sqrt{\Delta}}{2a}$ et $\dfrac{-b+\sqrt{\Delta}}{2a}$.
\item Si $\Delta =0$, alors pour tout complexe $z$, $az^2+bz+c = a\left(z+\dfrac{b}{2a}\right)^2$. L'unique racine de $P$ est $-\dfrac{b}{2a}$.
\item Si $\Delta <0$, alors $-\Delta >0$. Ainsi, pour tout complexe $z$,
\[az^2+bz+c = a\left(\left(z+\dfrac{b}{2a}\right)^2+ \left(- \dfrac{\Delta}{4a^2}\right)\right)= a\left(\left(z+\dfrac{b}{2a}\right)^2- i^2 \times \left(\sqrt{\dfrac{-\Delta}{4a^2}}\right)^2\right).\]

On se retrouve dans le même cas que dans la propriété précédente, et on procède donc de la même manière. Pour tout complexe $z$,
\[az^2+bz+c=a\left(\left(z+\dfrac{b}{2a}\right)^2-   \left(\dfrac{i\sqrt{-\Delta}}{2a}\right)^2\right) = a\left(z+\dfrac{b-i\sqrt{-\Delta}}{2a}\right)\left(z+\dfrac{b+i\sqrt{-\Delta}}{2a}\right).\]
Ainsi, $P$ admet deux racines complexes conjuguées : $\dfrac{-b-i\sqrt{-\Delta}}{2a}$ et $\dfrac{-b+i\sqrt{-\Delta}}{2a}$.
\end{itemize}
\end{demonstration}

\begin{example}
On cherche à résoudre l'équation $2z^2+3z+5=0$. 

Le discriminant du polynôme $2z^2+3z+5$ vaut $3^2-4\times 2  \times 5=-31<0$.

L'équation $2z^2+3z+5=0$ possède donc deux solutions : $z_1=\dfrac{-3+i\sqrt{31}}{4}$ et $z_2=\dfrac{-3-i\sqrt{31}}{4}$. 
\end{example}

\textbf{Remarque} : La démonstration met en évidence la factorisation de tout polynôme de degré 2 à coefficients réels, peu importe la nature de leurs racines (réelles ou imaginaires). 

Si $z_1$ et $z_2$ sont les deux racines du polynôme $az^2+bz+c$, alors pour tout complexe $z$, on a \[az^2+bz+c=a(z-z_1)(z-z_2).\]



\chapter{Exercices : Nombres complexes (1)}
 
\section*{Ensemble $\mathbb{C}$ des nombres complexes}



\begin{exercise}[topic=cpx01] Exprimer les nombres complexes suivants sous forme algébrique.

\begin{tabularx}{\linewidth}{XXX}
$1+2i-4i+7-3i^2$ & $(1+2i)(5-4i)$ & $\left(\dfrac{1}{3}+\dfrac{2}{5}i\right)\left(\dfrac{5}{2}+\dfrac{1}{4}i\right)$\\

$(5+3i)(10-6i)$ & $(2i)^{10}$ & $\left(\dfrac{\sqrt{2}}{2}+i \dfrac{\sqrt{2}}{2}\right)^8$

\end{tabularx}
\vspace{-1cm}\end{exercise}

\begin{solution}
On a 
\begin{itemize}
\item \(1+2i-4i+7-3i^2=8-2i+3=11-2i\).
\item \((1+2i)(5-4i)=1\times 5 + 1 \times (-4i) + 2i \times 5 +2i \times (-4i) = 5-4i+10i+8=13+6i\).
\item \(\left(\dfrac{1}{3}+\dfrac{2}{5}i\right)\left(\dfrac{5}{2}+\dfrac{1}{4}i\right)=\dfrac{1}{3} \times \dfrac{5}{2} + \dfrac{1}{3} \times \dfrac{1}{4}i+\dfrac{2}{5}i \times \dfrac{5}{2} +\dfrac{2}{5}i \times \dfrac{1}{4}i = \dfrac{5}{6}+\dfrac{1}{12}i+i-\dfrac{1}{10}=\dfrac{11}{15}+\dfrac{13}{12}i\).
\item \((5+3i)(10-6i)=5\times 10 + 5 \times(-6i)+3i \times 10 +3i \times (-6i) = 50 -30i+30i+18=68\).
\item \((2i)^{10}=2^{10} \times i^{10}=1024 \times (i^2)^5=1024 \times (-1)^5=-1024\).
\item Calculons d'abord \(\left(\dfrac{\sqrt{2}}{2}+i \dfrac{\sqrt{2}}{2}\right)^2\) en utilisant une identité remarquable.

On a \(\left(\dfrac{\sqrt{2}}{2}+i \dfrac{\sqrt{2}}{2}\right)^2 = \left(\dfrac{\sqrt{2}}{2}\right)^2 + 2 \times \dfrac{\sqrt{2}}{2} \times \dfrac{\sqrt{2}}{2}i + \left(\dfrac{\sqrt{2}}{2}i\right)^2=\dfrac{2}{4}+i-\dfrac{2}{4}=i \).

Ainsi, \(\left(\dfrac{\sqrt{2}}{2}+i \dfrac{\sqrt{2}}{2}\right)^8 = \left(\left(\dfrac{\sqrt{2}}{2}+i \dfrac{\sqrt{2}}{2}\right)^2\right)^4=i^4=i\times i \times i \times i = (-1) \times (-1) = 1\)
\end{itemize}
\end{solution}

\begin{exercise}[topic=cpx01]Résoudre les équations suivantes, d'inconnue $z\in\mathbb{C}$.

\renewcommand{\arraystretch}{2}
\begin{tabularx}{\linewidth}{XXX}
$z+2i-4=5i+3$ & $\dfrac{z-4i}{3}=\dfrac{1}{2}+5i$ & $(3+2i)z+i-5=1+5i$ \\
$(z-3+i)(2z-5i)=0$ & $\dfrac{z+1+i}{1-z}=0$ & $\dfrac{z-3i+1}{z-3}=\dfrac{1}{4}$

\end{tabularx}
\vspace{-1cm}\end{exercise}

\begin{solution}
\begin{itemize}
\item \(z+2i-4=5i+3 \Leftrightarrow z = 3+4+5i-2i \Leftrightarrow z = 7+3i\).
\item \(\dfrac{z-4i}{3}=\dfrac{1}{2}+5i \Leftrightarrow z = 3 \times \left(\dfrac{1}{2}+5i\right)+4i \Leftrightarrow z = \dfrac{3}{2}+19i\).
\item \((3+2i)z+i-5=1+5i \Leftrightarrow (3+2i)z = 1+5+5i-i \Leftrightarrow (3+2i)z = 6+4i \Leftrightarrow (3+2i)z = 2\times (3+2i) \Leftrightarrow z = 2\).
\item \((z-3+i)(2z-5i)=0 \Leftrightarrow z-3+i=0\) ou \(2z-5i=0 \Leftrightarrow z = -3-i\) ou \(z= \dfrac{5}{2}i\).
\item Soit \(z\in\mathbb{C}\setminus \{1\}\),
\(\dfrac{z+1+i}{1-z}=0 \Leftrightarrow z+1+i = 0 \Leftrightarrow z=-1-i\).
\item Soit \(z \in \mathbb{C}\setminus \{3\}\).

\(\dfrac{z-3i+1}{z-3}=\dfrac{1}{4} \Leftrightarrow 4(z-3i+1)=z-3 \Leftrightarrow 4z-12i+4=z-3 \Leftrightarrow 3z = -7+12i \Leftrightarrow z = -\dfrac{7}{3}+4i\).
\end{itemize}
\end{solution}

\begin{exercise}[topic=cpx01]Résoudre les système suivants, d'inconnues complexes $z_1$ et $z_2$.


\begin{tabularx}{\linewidth}{XX}$\left\{\begin{array}{rcl}
2z_1+3z_2 &=& 5-3i\\
z_1-3z_2&=&4+6i\end{array}\right.$ & $\left\{\begin{array}{rcl}
3z_1+4z_2 &=& 1+2i\\
z_1-2z_2&=&2+i\end{array}\right.$ \end{tabularx}
\vspace{-0.5cm}\end{exercise}

\begin{solution}\hspace{0pt}

\textbf{Premier système}

\(\left\{\begin{array}{rcl}
2z_1+3z_2 &=& 5-3i\\
z_1-3z_2&=&4+6i\end{array}\right.\)\(\Leftrightarrow \left\{\begin{array}{rcl}
2z_1+3z_2 &=& 5-3i\\
3z_1 &=&9+3i \end{array}\right. \Leftrightarrow \left\{\begin{array}{rcl}
2z_1+3z_2 &=& 5-3i\\
z_1 &=&3+i \end{array}\right.\Leftrightarrow \left\{\begin{array}{rcl}
2(3+i)+3z_2 &=& 5-3i\\
z_1 &=&3+i \end{array}\right.\)

\(\Leftrightarrow \left\{\begin{array}{rcl}
6+2i+3z_2 &=& 5-3i\\
z_1 &=&3+i \end{array}\right.\Leftrightarrow \left\{\begin{array}{rcl}
3z_2 &=& -1-5i\\
z_1 &=&3+i \end{array}\right.\Leftrightarrow \left\{\begin{array}{rcl}
z_1 &=& -3+i\\
z_2 &=& -\dfrac{1}{3}-\dfrac{5}{3}i \end{array}\right.\)

\newpage 
\textbf{Second système}

\(\left\{\begin{array}{rcl}
3z_1+4z_2 &=& 1+2i\\
z_1-2z_2&=&2+i\end{array}\right.\)\(\Leftrightarrow \left\{\begin{array}{rcl}
3z_1+4z_2 &=& 1+2i\\
3z_1-6z_2 &=&6+3i\end{array}\right.\Leftrightarrow \left\{\begin{array}{rcl}
3z_1+4z_2 &=& 1+2i\\
10z_2 &=&-5-i  \end{array}\right.\Leftrightarrow \left\{\begin{array}{rcl}
3z_1+4z_2 &=& 1+2i\\
z_2 &=&-\dfrac{1}{2}-\dfrac{i}{10} \end{array}\right.\)

\(\Leftrightarrow \left\{\begin{array}{rcl}
3z_1+4\left(-\dfrac{1}{2}-\dfrac{i}{10}\right) &=& 1+2i\\
z_2 &=&-\dfrac{1}{2}-\dfrac{i}{10} \end{array}\right.\Leftrightarrow \left\{\begin{array}{rcl}
3z_1-2-\dfrac{2i}{5} &=& 1+2i\\
z_2 &=&-\dfrac{1}{2}-\dfrac{i}{10} \end{array}\right.\Leftrightarrow \left\{\begin{array}{rcl}
3z_1 &=& 3+\dfrac{12i}{5}\\
z_2 &=&-\dfrac{1}{2}-\dfrac{i}{10} \end{array}\right.\)

\(\Leftrightarrow \left\{\begin{array}{rcl}
z_1 &=& 1+\dfrac{4i}{5}\\
z_2 &=&-\dfrac{1}{2}-\dfrac{i}{10} \end{array}\right.\)

\end{solution}


\begin{exercise}[topic=cpx01]Soit $x$ un réel et $z=2x+1+i(x^2+5x-4)$
\begin{enumerate}
\item Déterminer la valeurs de $x$ pour laquelle $z$ est un imaginaire pur. Que vaut alors $z$ ?
\item Existe-t-il des valeurs de $x$ pour lesquelles $z$ est réel ? Que vaut alors $z$ ?

\end{enumerate}\end{exercise}

\begin{solution}
\(z\) est imaginaire pur si sa partie réelle est nulle, c'est-à-dire \(2x+1=0\) et donc \(x=-\dfrac{1}{2}\). Dans ce cas,
\[z=\left( \left(-\dfrac{1}{2}\right)^2+5\times \left(-\dfrac{1}{2}\right)-4\right)i=\left(\dfrac{1}{4}-\dfrac{5}{2}-4\right)i=-\dfrac{25i}{4}.\]

\(z\) est réel si sa partie imaginaire est nulle, c'est-à-dire \(x^2+5x-4=0\). Les solutions sont \(-\dfrac{5}{2}-\dfrac{\sqrt{41}}{2}\) et \(-\dfrac{5}{2}+\dfrac{\sqrt{41}}{2}\). Dans le premier cas, on a alors \(z=-4-\sqrt{41}\). Dans le second cas, on a \(z=-4+\sqrt{41}\).

\end{solution}


\begin{exercise}[topic=cpx01]Soit $z$ et $z'$ deux complexes tels que $Re(zz')=Re(z) \times Re(z')$. Montrer que $z$ ou $z'$ est réel.\end{exercise}

\begin{solution}Notons $z=a+ib$ et $z'=a'+ib'$. On a alors $Re(z) \times Re(z')=aa'$.

Par ailleurs, $zz'=aa'-bb'+i(ab'+a'b)$ et donc $Re(zz')=aa'-bb'$. Ainsi, $Re(zz')=Re(z)\times Re(z')$ si et seulement si $aa'-bb'=aa'$ soit $bb'=0$ et donc $b=0$ ou $b'=0$. Autrement dit, $z$ ou $z'$ est réel.\end{solution}

\begin{exercise}[topic=cpx01]On considère la fonction $f$ définie pour tout $z$ dans $\mathbb{C}$ par $f(z)=z^2+2z+9$.
\begin{enumerate}
\item Soit $z=a+ib$ un complexe. Exprimer les parties réelles et imaginaires de $f(z)$ en fonction de $a$ et $b$.
\item Quels sont les nombres complexes dont l'image par $f$ est un réel ?
\end{enumerate}\end{exercise}

\begin{solution}

\(f(a+ib)=(a+ib)^2+2(a+ib)+9=a^2+2aib+(ib)^2+2a+2ib+9=a^2-b^2+2a+i \times 2b(a+1)\).

\(f(a+ib)\) est réel si sa partie imaginaire, c'est-à-dire \( 2b(a+1)\) vaut 0. Ainsi, \(f(z)\) est réel si et seulement si \(z\) est réel ou \(Re(z)=-1\).
\end{solution}


\begin{exercise}[topic=cpx01]Pour tout complexe $z$, on note $P(z)=6z^3+(1-30i)z^2-(12+5i)z+60i$.
\begin{enumerate}
\item Montrer que l'équation $P(z)=0$ admet une unique solution imaginaire pure que l'on déterminera.
\item Montrer qu'il existe trois réels $a$, $b$ et $c$ tels que, pour tout complexe $z$, $P(z)=(z-5i)(az^2+bz+c)$.
\item En déduire toutes les solutions de l'équation $P(z)=0$.
\end{enumerate}\end{exercise}

\begin{solution}Soit $b$ un réel. On a alors
\[P(ib)=6(ib)^3+(1-30i)(ib)^2-(12+5i)ib+60i=-6ib^3-b^2+30ib^2-12ib+5b+60i\]
\[=-b^2+5b+i(-6b^3+30b^2+12b+60).\]
Ainsi, $P(ib)=0$ si et seulement si $-b^2+5b=0$ et $-6b^3+30b^2+12b+60=0$.

La première équation donne $b=0$ ou $b=5$. Testons donc ces deux valeurs.

On a $P(0i)=P(0)=60i\neq 0$. Par ailleurs, 
\[P(5i)=6 \times (5i)^3+(1-30i) \times (5i)^2-(12+5i)\times 5i+60i=-750i-25+750i-60i+25+60i=0\]
Ainsi, l'équation $P(z)=0$ admet bien une unique solution imaginaire pure. Cette solution est $5i$.

Pour tout complexe $z$, on a $(z-5i)(az^2+bz+c)=az^3+(b-5ai)z^2+(c-5bi)z-5ci$.

Ainsi, en identifiant les coefficients, on a $P(z)=(z-5i)(az^2+bz+c)$ si et seulement si $a=6$, $b-5ai=1-30i$, $c-5bi=-12-5i$ et $-5ci=60i$.

On obtient donc $a=6$, $b=1$ et $c=-12$. Pour tout complexe $z$, on a $P(z)=(z-5i)(6z^2+z-12)$.

Or, $6z^2+z-12$ est un polynôme du second degré à coefficients réels. 

Son discriminant vaut $1^2-4\times 6 \times (-12)=289>0$. Ce polynôme admet donc deux racines réelles qui sont $\dfrac{-1-\sqrt{289}}{2\times 6}=-\dfrac{3}{2}$ et $\dfrac{-b+\sqrt{289}}{2\times 6}=\dfrac{4}{3}$.

Les solutions de l'équation $P(z)=0$ sont donc $5i$, $-\dfrac{3}{2}$ et $\dfrac{4}{3}$.
\end{solution}


\begin{exercise}[topic=cpx01]Résoudre l'équation $z^2=-3+4i$ sur $\mathbb{C}$. On pourra par exemple poser $z=a+ib$.\end{exercise}

\begin{solution}Posons $z=a+ib$, avec $a$ et $b$ réels. On a alors $z^2=a^2-b^2+2iab=-3+4i$.

Ainsi, en identifiant parties réelles et imaginaires, on a $a^2-b^2=-3$ et $2ab=4$.\\ En particulier, $a$ et $b$ sont non nuls et on a donc $b=\dfrac{2}{a}$.

En remplaçant $b$ dans la première équation, on obtient $a^2-\dfrac{4}{a^2}=-3$ et donc $\dfrac{a^4+3a^2-4}{a^2}=0$.

Posons $x=a^2$, on a alors $x^2+3x-4=0$, un polynôme du second degré dont les racines sont 1 et $-4$.

On a donc $a^2=1$ ou $a^2=-4$. Puisque $a$ est réel, cette dernière égalité est impossible.\\ Finalement, on a forcément $a^2=1$ et donc $a=1$ ou $a=-1$.

Si $a=1$, alors $b=\dfrac{2}{a}=2$. Si $a=-1$, alors $b=\dfrac{2}{a}=-2$. On a donc $z=1+2i$ ou $z=-1-2i$.

Vérifions :
\begin{itemize}
\item $(1+2i)^2=1^2+4i+(2i)^2=1+4i-4=-3+4i$.
\item $(-1-2i)^2=(-(1+2i))^2=-3+4i$.
\end{itemize}
Les solutions de l'équation $z^2=-3+4i$ sont donc $1+2i$ et $-1-2i$.
\end{solution}
\printcollection{cpx01}

\section*{Conjugué d'un nombre complexe}



\begin{exercise}[topic=cpx02]Soit $z=a+ib$ un nombre complexe. Exprimer le produit $z\overline{z}$ en fonction de $a$ et $b$. \end{exercise}

\begin{solution}On a \(z\overline{z}=(a+ib)(a-ib)=a^2-(ib)^2=a^2+b^2\). Une formule bien utile et qui nous sera bien utile dans un prochain chapitre !\end{solution}

\begin{exercise}[topic=cpx02]Démontrer la propriété du cours suivante : Soit $z\in\mathbb{C}$. On a $z\in i\mathbb{R}$ si et seulement si $\overline{z}=-z$.\end{exercise}

\begin{solution}Soit \(z=a+ib\) un complexe.

Si \(z\) est un imaginaire pur, alors \(a=0\). Ainsi, \(z=ib\) et \(\overline{z}=-ib=-z\).

Si \(z=-\overline{z}\), alors \(a+ib=-(a-ib)=-a+ib\). Il en résulte que \(2a=0\) et donc \(a=0\). \(z\) est donc imaginaire pur.\end{solution}

\begin{exercise}[topic=cpx02]Soit $z$ un nombre complexe.
\begin{enumerate}
\item Montrer que $Re(z)=\dfrac{z+\overline{z}}{2}$
\item Trouver et démontrer une relation similaire pour exprimer $Im(z)$ en fonction de $z$ et $\overline{z}$.
\end{enumerate}\end{exercise}

\begin{solution}
Soit \(z=a+ib\) un complexe.
\[\dfrac{z+\overline{z}}{2}=\dfrac{a+ib+a-ib}{2}=\dfrac{2a}{2}=a=Re(z).\]
De même,
\[\dfrac{z-\overline{z}}{2i}=\dfrac{a+ib-(a-ib)}{2}=\dfrac{2ib}{2i}=b=Im(z).\]
\end{solution}

\begin{exercise}[topic=cpx02]On considère le complexe $z=(1-3i)(5+4i)(1+3i)(10-8i)$. Sans calcul, montrer que $z\in\mathbb{R}$.\end{exercise}

\begin{solution}D'une part, remarquons que \(10-8i=2(5-4i)\). Ainsi, \(z=2(1-3i)(5+4i)(1+3i)(5-4i)\)

D'après les formules sur le produit de conjugués, on a
\[\overline{z}=\overline{2(1-3i)(5+4i)(1+3i)(5-4i)}=\overline{2}\times\overline{(1-3i)} \times \overline{(5+4i)} \times \overline{(1+3i)} \times \overline{(5-4i)}.\]

Ainsi,
\[\overline{z}=2(1+3i)(5-4i)(1-3i)(5+4i)=z\]
Puisque \(\overline{z}=z\), on en conclut que \(z\) est réel.\end{solution}


\begin{exercise}[topic=cpx02]Résoudre les équations suivantes, d'inconnue $z\in\mathbb{C}$.

\begin{tabularx}{\linewidth}{XXX}
$2z+3=3\overline{z}+4-i$ & $2z-4\overline{z}= 6-3i$ & $z+\overline{z}=3+4i$ \\

$2z+4+i=(3+i)\overline{z}+2-3i$  &  $\overline{z}-1=z\overline{z}-i$ & $\overline{z}=iz$

\end{tabularx}\end{exercise}

\begin{solution}
\begin{itemize}
\item Soit \(z=a+ib\in\mathbb{C}\). 

On a \(2z+3=3\overline{z}+4-i \Leftrightarrow 2(a+ib)+3=3(a-ib)+4-i \Leftrightarrow 2a-3a+3-4+i(2b+3b+1)=0 \Leftrightarrow -a-1 +i(5b+1)=0\).

On a donc le système \(\left\{\begin{array}{rcl}-a-1&=&0 \\ 5b+1&=&0\end{array}\right.\) ce qui conduit à \(a=-1\) et \(b=-\dfrac{1}{5}\)

La solution de cette équation est donc \(-1-\dfrac{i}{5}\).
\vskip5pt
\item Soit \(z=a+ib\in\mathbb{C}\). 

On a \(2z-4\overline{z}= 6-3i \Leftrightarrow 2(a+ib)-4(a-ib)=6-3i \Leftrightarrow 2a-4a-6+i(2b+4b+3) = 0 \Leftrightarrow -2a-6+i(6b+3)=0\).

On a donc le système \(\left\{\begin{array}{rcl}-2a-6&=&0 \\ 6b+3 &=&0\end{array}\right.\) ce qui conduit à \(a=-3\) et \(b=-\dfrac{1}{2}\).

La solution de cette équation est donc \(-3-\dfrac{i}{2}\).
\vskip5pt
\item Soit \(z=a+ib\in\mathbb{C}\).  

On a \(z+\overline{z}=3+4i \Leftrightarrow a+ib+a-ib=3+4i \Leftrightarrow 2a-3-4i = 0\).

Or, la partie imaginaire du membre de gauche est non nulle. Cette équation n'admet aucune solution (on aurait également tout simplement pu utiliser un exercice précédent : \(z+\overline{z} = 2Re(z)\).
\vskip5pt
\item Soit \(z=a+ib\in\mathbb{C}\). 

On a \(2z+4+i=(3+i)\overline{z}+2-3i \Leftrightarrow 2(a+ib)+4+i=(3+i)(a-ib)+2-3i \Leftrightarrow 2a+2ib+i+4=3a-3ib+ia+b+2-3i\).

\(\Leftrightarrow 2a-3a -b-2+4 +i(2b+1+3b-a+3) \Leftrightarrow -a-b+2+i(-a+5b+4)=0\).

Ceci nous amène à résoudre le système 

\(\left\{\begin{array}{rcl}-a-b+2&=&0 \\ -a+5b+4 &=&0\end{array}\right. \Leftrightarrow \left\{\begin{array}{rcl}-a-b+2&=&0 \\ -6b-2 &=&0\end{array}\right. \Leftrightarrow \left\{\begin{array}{rcl}-a-b+2&=&0 \\ b &=& - \dfrac{1}{3} \end{array}\right.\)

\( \Leftrightarrow \left\{\begin{array}{rcl}-a-\left(-\dfrac{1}{3}\right)+2&=&0 \\ b &=& - \dfrac{1}{3} \end{array}\right.\Leftrightarrow \left\{\begin{array}{rcl} a &=& \dfrac{7}{3} \\ b &=& - \dfrac{1}{3} \end{array}\right.\)

La solution de cette équation est \(\dfrac{7}{3}-\dfrac{i}{3}\).
\vskip5pt
\item Soit \(z=a+ib\in\mathbb{C}\). 

On a  \(\overline{z}-1=z\overline{z}-i \Leftrightarrow a-ib-1=(a+ib)(a-ib)-i \Leftrightarrow a-ib-1=a^2+b^2-i \Leftrightarrow a-a^2-b^2 +i(b+1)=0\).

On a donc le système  \(\left\{\begin{array}{rcl}a-a^2-b^2&=&0 \\ b+1 &=&0\end{array}\right.\).

La seconde ligne nous donne \(b=-1\). En remplaçant \(b\) par \(-1\) dans la première équation, on a alors \(a-a^2-1=0\). C'est un polynôme du second degré dont le discriminant vaut \(1^2-4\times(-1) \times (-1) = -3\). Ainsi, l'équation \(a-a^2-1=0\) n'admet pas de solution réelle. L'équation de départ n'a donc pas de solution.
\vskip5pt
\item Soit \(z=a+ib\in\mathbb{C}\) 

On a \(\overline{z}=iz \Leftrightarrow a-ib = i(a+ib) \Leftrightarrow a+b -i(a+b)=0\).

Cette équation est équivalente à \(a=-b\). Les solutions sont donc les nombres complexes \(a-ia\) pour tout réel \(a\).
\end{itemize}
\end{solution}


\begin{exercise}[topic=cpx02]Résoudre dans $\mathbb{C}$ l'équation $z^2-\overline{z}+1=0$.\end{exercise}

\begin{solution}Soit \(z=a+ib\in\mathbb{C}\). On a 
 \[z^2-\overline{z}+1=0 \Leftrightarrow (a+ib)^2-(a-ib)+1=0 \Leftrightarrow a^2 +2aib -b^2-a+ib+1=0 \Leftrightarrow a^2-b^2-a+1+ib(2a+1)=0.\]

On a donc le système \(\left\{\begin{array}{rcl}a^2-b^2-a+1&=&0 \\ b(2a+1) &=&0\end{array}\right.\).

La deuxième ligne nous amène à différencier deux cas :
\begin{itemize}
\item Soit \(b=0\). La première équation devient alors \(a^2-a+1=0\). C'est un polynôme du second degré dont le discriminant vaut \(-3\). Il ne possède donc pas de racine réelle.
\item Soit \(a=-\dfrac{1}{2}\). La première équation devient alors \( \left(-\dfrac{1}{2}\right)^2+\dfrac{1}{2}+1=b^2\) soit \(b^2=\dfrac{7}{4}\) ce qui conduit à \(b=\dfrac{\sqrt{7}}{2}\) ou \(b=-\dfrac{\sqrt{7}}{2}\).\end{itemize}

Les solutions sont donc \(-\dfrac{1}{2}+\dfrac{i\sqrt{7}}{2}\) et \(-\dfrac{1}{2}-\dfrac{i\sqrt{7}}{2}\). \end{solution}


\printcollection{cpx02}

\section*{Division dans $\mathbb{C}$}

\begin{exercise}[topic=cpx03]Ecrire les nombres suivants sous forme algébrique.

\begin{tabularx}{\linewidth}{XXXXXX}
$\dfrac{1}{i}$ & $\dfrac{1}{1+i}$ & $\dfrac{2}{4-5i}$ & $\dfrac{2i}{i-3}$ & $\dfrac{1-6i}{2+4i}$ & $\dfrac{3-2i}{2i-3}$
\end{tabularx}\end{exercise}

\begin{solution}On a ...
\begin{itemize}
\item \( \dfrac{1}{i} = \dfrac{-i}{i \times -i} = \dfrac{-i}{1}=-i\).
\vskip5pt
\item \(\dfrac{1}{1+i} = \dfrac{1-i}{(1+i)(1-i)}=\dfrac{1-i}{1^2-i^2}=\dfrac{1-i}{2}=\dfrac{1}{2}-\dfrac{i}{2}\).
\vskip5pt
\item \(\dfrac{2}{4-5i} = \dfrac{2(4+5i)}{(4-5i)(4+5i)}=\dfrac{8+10i}{4^2-(5i)^2}=\dfrac{8}{41}+\dfrac{10i}{41}\).
\vskip5pt
\item \(\dfrac{2i}{i-3} = \dfrac{(2i)(-i-3)}{(i-3)(-i-3)}=\dfrac{2-6i}{3^2-i^2}=\dfrac{2}{10}-\dfrac{6i}{10}=\dfrac{1}{5}-\dfrac{3i}{5}\).
\vskip5pt
\item \(\dfrac{1-6i}{2+4i} = \dfrac{(1-6i)(2-4i)}{(2+4i)(2-4i)}=\dfrac{2-4i-12i-24}{2^2-(4i)^2}=-\dfrac{22}{20}-\dfrac{16i}{22}=-\dfrac{11}{10}-\dfrac{8i}{11}\).
\vskip5pt
\item \(\dfrac{3-2i}{2i-3}=\dfrac{-(2i-3)}{2i-3}=-1\). Attention à ne pas se lancer précipitamment dans les calculs !
\end{itemize}\end{solution}

\begin{exercise}[topic=cpx03]Résoudre les équations suivantes, d'inconnues complexes $z$. On exprimera le résultat sous forme algébrique.

\begin{tabularx}{\linewidth}{XXX}
$(1+2i)z+3=5-4i$ & $(3-4i)z+1-3i=2iz$ & $\dfrac{z-i}{iz-1}=\dfrac{1}{3}$
\end{tabularx}

\end{exercise}

\begin{solution}
\begin{itemize}
\item Soit \(z\in\mathbb{C}\). On a \((1+2i)z+3=5-4i \Leftrightarrow (1-2i)z=2-4i \Leftrightarrow z=\dfrac{2-4i}{1+2i}\)

On met alors \(\dfrac{2-4i}{1+2i}\) sous forme algébrique : \(\dfrac{2-4i}{1+2i}=\dfrac{(2-4i)(1-2i)}{(1+2i)(1-2i)}=\dfrac{2-4i-4i-8}{1^2-(2i)^2}=-\dfrac{6}{5}-\dfrac{8i}{5}\)

La solution de cette équation est donc \(-\dfrac{6}{5}-\dfrac{8i}{5}\).

\item Soit \(z\in\mathbb{C}\). 

On a \((3-4i)z+1-3i=2iz \Leftrightarrow (3-4i)z-2iz=-1+3i \Leftrightarrow (3-6i)z=-1+3i\Leftrightarrow z=\dfrac{-1+3i}{3-6i}\).

On met alors \(\dfrac{-1+3i}{3-6i}\) sous forme algébrique.

\(\dfrac{-1+3i}{3-6i}=\dfrac{(-1+3i)(3+6i)}{(3-6i)(3+6i)}=\dfrac{-3-6i+9i-18}{3^2-(6i)^2}=-\dfrac{21}{45}+\dfrac{3i}{45}=-\dfrac{7}{15}+\dfrac{i}{15}\).

La solution de cette équation est donc \(-\dfrac{7}{15}+\dfrac{i}{15}\).

\item Voyons d'abord quand le dénominateur s'annule. Soit \(z\in\mathbb{C}\). On a \(iz-1=0 \Leftrightarrow iz=1 \Leftrightarrow z=\dfrac{1}{i} \Leftrightarrow z = -i\).

Soit donc \(z \in \mathbb{C}\setminus \{-i\}\).
\[\dfrac{z-i}{iz-1}=\dfrac{1}{3} \Leftrightarrow 3(z-i)=iz-1 \Leftrightarrow 3z-iz=-1+3i \Leftrightarrow (3-i)z=(-1+3i) \Leftrightarrow z =\dfrac{-1+3i}{3-i}.\]

On met alors \(\dfrac{-1+3i}{3-i}\) sous forme algébrique.
\[\dfrac{-1+3i}{3-i}=\dfrac{(-1+3i)(3+i)}{(3-i)(3+i)}=\dfrac{-3-i+9i-3}{3^2-i^2}=-\dfrac{6}{45}+\dfrac{8i}{45}=-\dfrac{3}{5}+\dfrac{4i}{5}.\]
La solution de cette équation est donc \(-\dfrac{3}{5}+\dfrac{4i}{5}\).\end{itemize}
\end{solution}


\begin{exercise}[topic=cpx03]Pour tout complexe $z$ différent de 1, on pose $f(z)=\dfrac{2-iz}{1-z}$
\begin{enumerate}
\item Montrer que $f(z)$ ne peut pas être égal à $i$.
\item Soit $Z$ un complexe différent de $i$. Déterminer, s'il(s) existe(nt), le(s) antécédent(s) de $Z$ par $f$.
\end{enumerate}

\end{exercise}

\begin{solution}Supposons qu'il existe un complexe \(z\) différent de 1 tel que \(\dfrac{2-iz}{1-z}=i\). 

On a alors \(2-iz=i(1-z)\) et donc \(2-iz=i-iz\) et finalement \(2=i\). C'est absurde ! Ainsi, pour tout complexe \(z\) différent de 1, \(f(z)\) est différent de \(i\).

Soit \(Z\) un complexe différent de \(i\) et \(z\) un complexe différent de 1.

\(f(z)=Z \Leftrightarrow \dfrac{2-iz}{1-z}=Z \Leftrightarrow 2-iz = Z(1-z) \Leftrightarrow zZ-iz = Z - 2 \Leftrightarrow z(Z-i)=Z-2 \Leftrightarrow z = \dfrac{Z-2}{Z-i}\).\end{solution}


\begin{exercise}[topic=cpx03]Soit $z$ un complexe non nul. Montrer de deux manières différentes que $\dfrac{1}{z}+\dfrac{1}{\overline{z}}$ est un réel.\newpage \end{exercise}

\begin{solution}

\textbf{Méthode 1 :} D'après le cours, on sait que \(\overline{\left(\dfrac{1}{z}\right)} = \dfrac{1}{\overline{z}} \). Or, on a vu plus tôt dans un exercice que pour tout complexe \(Z\), \(Z+\overline{Z}=Re(Z)\). En particulier, \(Z+\overline{Z}\) est réel. \(\dfrac{1}{z}+\dfrac{1}{\overline{z}}\) est donc un réel.

\textbf{Méthode 2 :} Soit \(z=a+ib\) un complexe,
\[\dfrac{1}{z}+\dfrac{1}{\overline{z}} = \dfrac{1}{a+ib}+\dfrac{a-ib}=\dfrac{a-ib}{(a+ib)(a-ib)}+\dfrac{a+ib}{(a+ib)(a-ib)}=\dfrac{2a}{a^2+b^2}.\]

Ainsi, \(a\) et \(b\) étant réels, \(\dfrac{1}{z}+\dfrac{1}{\overline{z}}\) est aussi un réel.\end{solution}


\begin{exercise}[topic=cpx03]Soit $j=-\dfrac{1}{2}+\dfrac{i\sqrt{3}}{2}$. Montrer que $j^2=\dfrac{1}{j}=\overline{j}$.\end{exercise}

\begin{solution}D'une part,
\[j^2=\left(-\dfrac{1}{2}+\dfrac{i\sqrt{3}}{2}\right)^2=\left(\dfrac{1}{2}\right)^2+2\times \left(-\dfrac{1}{2}\right) \times \dfrac{\sqrt{3}}{2} + \left(\dfrac{\sqrt{3}}{2}\right)^2=\dfrac{1}{4}-\dfrac{i\sqrt{3}}{2}-\dfrac{3}{4}=-\dfrac{1}{2}-\dfrac{i\sqrt{3}}{2}\]

Ainsi, \(j^2=\overline{j}\). De plus, 
\[j^3 = j^2 \times j = \left(-\dfrac{1}{2}-\dfrac{i\sqrt{3}}{2}\right) \left(-\dfrac{1}{2}+\dfrac{i\sqrt{3}}{2}\right)=\left(\dfrac{1}{2}\right)^2-\left(\dfrac{i\sqrt{3}}{2}\right)^2=\dfrac{1}{4}+\dfrac{3}{4}=1\]

On a donc \(j^3=1\), soit \(j^2 \times j =1\) et donc \(j^2 = \dfrac{1}{j}\).\end{solution}

\begin{exercise}[topic=cpx03]On considère la suite $(z_n)$ définie par $z_0=1$ et, pour tout entier naturel $n$, $z_{n+1}=\dfrac{z_n-6}{1+i}$
\begin{enumerate}
\item Exprimer $z_1$ et $z_2$ sous forme algébrique.
\item Pour tout entier naturel $n$, on pose $u_n=z_n-6i$
\begin{enumerate}
\item Exprimer $u_{n+1}$ en fonction de $u_n$. Quelle est la nature de la suite $(u_n)$ ?
\item En déduire une expression de $u_n$ puis de $z_n$ en fonction de $n$.
\end{enumerate}
\item On considère la suite $(t_n)$ définie par $t_0=1$ et, pour tout entier naturel $n$, $t_{n+1}=\dfrac{t_n-6}{1-i}$. 
\begin{enumerate}
\item Montrer par récurrence que pour tout entier naturel $n$, $t_n = \overline{z_n}$
\item Sans avoir recours à un nouveau calcul, exprimer $t_n$ en fonction de $n$ pour tout entier naturel $n$.
\end{enumerate}\end{enumerate}\end{exercise}

\begin{solution}\hspace{0pt}

\begin{enumerate}\item 
On a \begin{itemize}\item \(z_1=\dfrac{1-6}{1+i}=\dfrac{-5(1-i)}{(1+i)(1-i)}=\dfrac{-5+5i}{1^2-i^2}=-\dfrac{5}{2}+\dfrac{5i}{2}\).
\item \(z_2 = \dfrac{-\frac{5}{2}+\frac{5i}{2}+1}{1+i}=\dfrac{\left(-\frac{3}{2}+\dfrac{5i}{2}\right)(1-i)}{(1+i)(1-i)}=\dfrac{1}{2} \times \left(-\dfrac{3}{2}+\dfrac{3i}{2}+\dfrac{5i}{2}+\dfrac{5}{2}\right)=\dfrac{1+4i}{2}\).\end{itemize}

\item 
\begin{enumerate} 
\item Pour tout entier naturel \(n\), 
\[ u_{n+1}=z_{n+1}-6i = \dfrac{z_n-6}{1+i} - 6i = \dfrac{z_n-6-6i(1+i)}{1+i}=\dfrac{z_n-6-6i+6}{1+i}=\dfrac{z_n-6i}{1+i}=\dfrac{u_n}{1+i}\]
La suite \((u_n)\) est géométrique de raison \(\dfrac{1}{1+i}\).

\item On sait que \(u_0 = z_0 -6i = 1-6i\). Pour tout entier naturel \(n\), on a donc \(u_n = \dfrac{1-6i}{(1+i)^n}\) et donc \(z_n=u_n+6i=\dfrac{1-6i}{(1+i)^n}+6i\).\end{enumerate}
\item \begin{enumerate}
\item Pour tout entier naturel \(n\), on considère la proposition \(P(n)\) : \(t_n = \overline{z_n}\).
\begin{itemize}\item \textbf{Initialisation} : On a \(t_0=1\) et \(z_0=1\). Ainsi, \(t_0 = \overline{z_0}\). \(P(0)\) est vraie.
\item \textbf{Hérédité} : Soit \(n\in \mathbb{N}\). Supposons que \(P(n)\) est vraie. On a donc \(t_n = \overline{z_n}\). Or,
\[t_{n+1}=\dfrac{t_n-6}{1-i}=\dfrac{\overline{z_n}-6}{1-i}=\dfrac{\overline{z_n-6}}{\overline{1+i}}=\overline{\left(\dfrac{z_n-6}{1+i}\right)}=\overline{z_{n+1}}\]
Ainsi, \(P(n+1)\) est vraie.
\item \textbf{Conclusion} : \(P(0)\) est vraie et \(P\) est héréditaire. Par récurrence, \(P(n)\) est vraie pour tout entier naturel \(n\).\end{itemize}
\item Pour tout entier naturel \(n\),
\[t_n = \overline{z_n}=\overline{\dfrac{1-6i}{(1+i)^n}+6i}=\dfrac{1+6i}{(1-i)^n}-6i.\]\end{enumerate}\end{enumerate}
\end{solution}

\begin{exercise}[topic=cpx03]On considère la fonction $f:z\mapsto \dfrac{z-i}{z+i}$, définie sur $\mathbb{C} \setminus \{-i\}$.

\begin{enumerate}
\item Exprimer $f(1+2i)$ sous forme algébrique.
\item Montrer que pour tout complexe $z \neq -i$, $f(z) \neq 1$.
\item Soit $Z$ un complexe différent de 1. Déterminer l'unique antécédent de $Z$ par $f$.
\item Soit $x$ un réel. Montrer que $f(x) \times \overline{f(x)}=1$.
\item Réciproquement, soit $z$ un complexe différent de $-i$ tel que $f(z) \times \overline{f(z)}=1$. Montrer que $z$ est réel.
\end{enumerate}\end{exercise}

\begin{solution}On a \(f(1+2i)=\dfrac{1+2i-i}{1+2i+i}=\dfrac{1+i}{1+3i}=\dfrac{(1+i)(1-3i)}{(1+3i)(1-3i)}=\dfrac{1-3i+i+3}{1^2-(3i)^2}=\dfrac{4}{10}-\dfrac{2i}{10}=\dfrac{2}{5}-\dfrac{i}{5}\).

Supposons qu'il existe un complexe \(z\) différent de \(-i\) tel que \(f(z)=1\).

On a alors  \(\dfrac{z-i}{z+i} = 1\) et donc \(z-i=z+i\) et finalement \(2i=0\), ce qui est absurde. Ainsi, pour tout complexe \(z \neq -i\), \(f(z) \neq 1\).

Soit \(Z\) un complexe différent de 1 et \(z\) un complexe différent de \(-i\),
\[f(z)=Z \Leftrightarrow \dfrac{z-i}{z+i} = Z \Leftrightarrow z-i=Z(z+i)\Leftrightarrow z-zZ=i+Zi \Leftrightarrow z(1-Z)=i+Zi \Leftrightarrow z = \dfrac{i(1+Z)}{1-Z}.\]

Soit \(x\) un réel. On a alors \(\overline{x}=x\). Ainsi, \(f(x) \times \overline{f(x)} = \dfrac{x-i}{x+i} \times \overline{\left(\dfrac{x-i}{x+i}\right)} = \dfrac{x-i}{x+i} \times \dfrac{x+i}{x-i}=1.\)

Soit \(z\) un complexe différent de \(-i\) tel que \(f(z) \times \overline{f(z)}=1\). 

On a alors \(\dfrac{z-i}{z+i} \times \overline{\left(\dfrac{z-i}{z+i}\right)}=1 \) et donc \( \dfrac{z-i}{z+i} \times \dfrac{\overline{z}+i}{\overline{z}-i}=1 \).

En développant, on obtient que \(\dfrac{z\overline{z}-i\overline{z}+iz+1}{z\overline{z}+i\overline{z}-iz+1}=1\) et donc \(z\overline{z}-i\overline{z}+iz+1 = z\overline{z}+i\overline{z}-iz+1.\)

On en déduit que \(2i(z-\overline{z})=0.\) Or, \(z-\overline{z}=\dfrac{Im(z)}{2i}\). Ainsi
\(Im(z)=0.\)
\(z\) est donc un réel.\end{solution}

\printcollection{cpx03}

\section*{Équations du second degré}

\begin{exercise}[topic=cpx04]Résoudre les équations suivantes, d'inconnue complexe $z$.

\begin{tabularx}{\linewidth}{XXX}
$z^2-2z+5=0$ & $z^2=-9$ & $2z^2=3-\pi$ \\
$4z^2+3z+2=0$ & $z^3+2z^2+3z=0$ & $(z^2+3)(z^2+z+1)=0$.
\end{tabularx}\end{exercise}

\begin{solution}\hspace{0pt}

\begin{itemize}
\item Le discriminant de $z^2-2z+5$ vaut $(-2)^2-4 \times 5=-16<0$. L'équation $z^2-2z+5=0$ admet donc deux solutions complexes conjuguées qui sont $z_1=\dfrac{-(-2)+i\sqrt{16}}{2}=1+2i$ et $z_2=1-2i$.
\vskip5pt
\item Les solutions de $z^2=-9$ sont $z_1=3i$ et $z_2=-3i$.
\vskip5pt
\item Remarquons que $3-\pi<0$. On a $2z^2=3-\pi$ si et seulement si $z^2=\dfrac{3-\pi}{2}$ et donc $z=i\sqrt{\dfrac{\pi-3}{2}}$ ou $z=-i\sqrt{\dfrac{\pi-3}{2}}$.
\vskip5pt
\item Le discriminant de $4z^2+3z+2$ vaut $3^2-4 \times 4 \times 2 =23$. L'équation $4z^2+3z+2=0$ admet donc deux solutions complexes conjuguées qui sont $z_1=\dfrac{-3+i\sqrt{23}}{8}$ et $z_2=\dfrac{-3-i\sqrt{23}}{8}$.
\vskip5pt
\item On a $z^3+2z^2+3z=0$ si et seulement si $z(z^2+2z+3)=0$, c'est-à-dire $z=0$ ou $z^2+2z+3=0$. Le discriminant de $s^2+2z+3$ vaut $2^2-4 \times 3 \times 1 =-8<0$. L'équation $z^2+2z+3=0$ admet donc deux solutions complexes conjuguées qui sont $z_1=\dfrac{-2-i\sqrt{8}}{2\times 1}=-1-i\sqrt{2}$ et $z_2=-1+i\sqrt{2}$. Les solutions de l'équation $z^3+2z^2+3z=0$ sont donc 0, $-1-i\sqrt{2}$ et $-1+i\sqrt{2}$.
\vskip5pt
\item On a $z^2+3=0$ si et seulement si $z^2=-3$ soit $z=i\sqrt{3}$ ou $z=-i\sqrt{3}$. Par ailleurs, $z^2+z+1$ est un polynôme du second degré dont le discriminant vaut $1^2-4 \times 1 \times 1 =-3<0$. Ce polynôme admet donc deux racines complexes conjuguées qui sont $\dfrac{-1-i\sqrt{3}}{2}$ et $\dfrac{-1+i\sqrt{3}}{2}$.

Les solutions de l'équation sont donc $-i\sqrt{3}$, $i\sqrt{3}$, $\dfrac{-1+i\sqrt{3}}{2}$ et $\dfrac{-1-i\sqrt{3}}{2}$.
\end{itemize}
\end{solution}


\begin{exercise}[topic=cpx04]Trouver deux nombres complexes dont la somme vaut 5 et le produit vaut 12.\end{exercise}

\begin{solution}Notons $z_1$ et $z_2$ les deux nombres recherchés. On a alors $z_1+z_2=5$ et donc $z_1=5-z_2$. De plus, on a $z_1z_2=12$ soit $(5-z_2)z_2=12$ et donc $-z_2^2+5z_2-12=0$. Il s'agit d'une équation du second degré. Le discriminant du polynôme $-z^2+5z-12$ vaut $5^2-4 \times (-12) \times (-1) = -23<0$. Ce polynôme admet donc deux racines complexes conjuguées $\dfrac{-5-i\sqrt{23}}{-2}$ et $\dfrac{-5+i\sqrt{23}}{-2}$. Ces deux racines sont les nombres recherchées.\end{solution}

\begin{exercise}[topic=cpx04]En utilisant un changement de variable, résoudre l'équation $z^4+2z^2-8=0$\end{exercise}

\begin{solution}
Posons $Z=z^2$. On a $z^4+2z^2-8=0$ si et seulement si $Z^2+2Z-8=0$. Il s'agit d'une équation du second degré dont les solutions sont réelles et valent $-4$ et 2. 

Ainsi, $z^4+2z^2-8=0$ si et seulement si $z^2=-4$ ou $z^2=2$. Les solutions de cette équation sont donc $2i$, $-2i$, $\sqrt{2}$ et $-\sqrt{2}$.
\end{solution}

\begin{exercise}[topic=cpx04]Pour tout complexe $z$, on pose $P(z)=z^3-3z^2+7z-5$.
\begin{enumerate}
\item Déterminer une solution « évidente » de l'équation $P(z)=0$. On note $\alpha$ cette solution.
\item Déterminer des réels $a$, $b$ et $c$ tels que pour tout complexe $z$, $P(z)=(z-\alpha)(az^2+bz+c)$.
\item Résoudre l'équation $P(z)=0$ sur $\mathbb{C}$.
\end{enumerate}\end{exercise}

\begin{solution}
Remarquons que $P(1)=0$.

On a alors $(z-1)(az^2+bz+c)=az^3+(b-a)z^2+(c-b)z-c$. Ainsi, $P(z)=(z-1)(az^2+bz+c)$ si et seulement si $a=1$, $b-a=-3$, $c-b=7$ et $-c=-5$ soit  $a=1$, $b=-2$ et $c=5$.

Ainsi, pour tout complexe $z$, $P(z)=(z-1)(z^2-2z+5)$.

$z^2-2z+5$ est un polynôme du second degré dont le discriminant vaut $(-2)^2-4\times 1 \times 5=-16<0$. Ses racines sont donc $\dfrac{-(-2)+i\sqrt{16}}{2 \times 1}=1+2i$ et $1-2i$.

Les solutions de l'équation $P(z)=0$ sont donc $1$, $1+2i$ et $1-2i$.
\end{solution}

\begin{exercise}[topic=cpx04]Soit $x\in]0;\pi[$. Résoudre l'équation $z^2+2z\cos(x)+1=0$, d'inconnue complexe $z$.\end{exercise}

\begin{solution}Le discriminant de $z^2+2z\cos(x)+1$ vaut $(2\cos(x))^2-4\times 1 \times 1=4(cos^2(x)-1)=-4\sin^2(x)<0$ puisque $x\in]0;\pi[$. De plus, sur cet intervalle, on a $\sin(x)>0$.

Les racines de ce polynôme sont donc $\dfrac{-2\cos(x)+i\sqrt{4\sin^2(x)}}{2 \times 1}=-\cos(x)+i\sin(x)$ et $-\cos(x)-i\sin(x)$.\end{solution}

\printcollection{cpx04}

\chapter{Corrigés}

\section*{Ensemble $\mathbb{C}$ des nombres complexes}

\printsolutions[collection={cpx01}, headings={false} ]

\section*{Conjugué d'un nombre complexe}

\printsolutions[collection={cpx02}, headings={false} ]

\section*{Division dans $\mathbb{C}$}

\printsolutions[collection={cpx03}, headings={false} ]

\section*{Équations du second degré}

\printsolutions[collection={cpx04}, headings={false} ]

\end{document}